% !PRINTED 30
\weberpage{30}
\begin{center}
\S\,8. \textbf{\gesperrt{The Jacobi Transformation Principle.}}
\end{center}

The \gesperrt{Jacobi} theory of transformation sets itself in general the task of
reducing an elliptic differential, by an algebraic substitution, to
another elliptic differential. In setting out here the foundations of
this theory, we make use of the representation of the elliptic
differential in homogeneous variables [\S\,1, (4)]:
\weq{1}{\frac{x\,dy - y\,dx}{\sqrt{f(x,y)}}\,,}
and substitute in it, for $x,y$, two coprime integral homogeneous
functions of the same, but arbitrary, degree $n$ in two new variables
$\xi,\eta$:
\weq{2}{x = U(\xi,\eta), \qquad y = V(\xi,\eta).}
From the equations
\[
\begin{aligned}
nU &= \xi\,\frac{\partial U}{\partial \xi} + \eta\,\frac{\partial U}{\partial \eta}, &
dU &= d\xi\,\frac{\partial U}{\partial \xi} + d\eta\,\frac{\partial U}{\partial \eta},\\
nV &= \xi\,\frac{\partial V}{\partial \xi} + \eta\,\frac{\partial V}{\partial \eta}, &
dV &= d\xi\,\frac{\partial V}{\partial \xi} + d\xi\,\frac{\partial V}{\partial \eta} % [sic] second differential printed d\xi; the total differential requires d\eta
\end{aligned}
\]
there then follows
\weq{3}{U\,dV - V\,dU = H(\xi\,d\eta - \eta\,d\xi),}
if by $H$ the functional determinant
\weq{4}{H = \frac{1}{n}\left(\frac{\partial U}{\partial \xi}\frac{\partial V}{\partial \eta} - \frac{\partial V}{\partial \xi}\frac{\partial U}{\partial \eta}\right)}
is meant.

Accordingly the elliptic differential (1) passes over into the following:
\weq{5}{H\,\frac{\xi\,d\eta - \eta\,d\xi}{\sqrt{f(U,V)}}\,.}
Now in order that this differential should again receive the form of an
elliptic one, it is necessary that from the function $f(U,V)$, whose
degree is the $4n$th, a square factor $T^2$ of degree $4n-4$ can be
separated off; that is, if $\varphi(\xi,\eta)$ denotes a function of the
fourth degree, that
\weq{6}{f(U,V) = T^2\,\varphi(\xi,\eta)}
should hold, whereby, if
\weq{7}{\frac{T}{H} = M}
%CONT>>
